\section{Minimal singularities of nef klt adjoints}
\label{met:section}

We prove the metric statement needed in the remaining nonprojective
case.  The proof does not require sections of a positive twist.  We use
\(\ddc=\ii\partial\bar\partial\), and identify the class of a line bundle
with its curvature class; thus the usual factor \(2\pi\) in its first
Chern class is understood.  Weights on a rational line bundle mean
weights obtained by taking a root of a metric on a fixed Cartier power.
All the line identifications below are actual rational line-bundle
identifications.

\begin{lemma}[Minimal metric of a nef klt adjoint]
\label{met:zero-lelong}
Let \(p\colon M\to Z\) be a compact K\"ahler log resolution of a normal
compact K\"ahler klt pair \((Z,\Delta)\), where \(\Delta\geq0\) is rational
and \(K_Z+\Delta\) is a nef rational line bundle.  Then a semipositive
metric with minimal singularities on
\[
 L=p^*(K_Z+\Delta)
\]
has zero Lelong numbers at every point of \(M\).
\end{lemma}

\begin{proof}
We may suppose that \(M\) is connected.  The assertion is immediate in
dimension zero, so write \(n=\dim M>0\).

The proof is by contradiction from a positive Lelong number. We construct
the normalized Monge--Amp\`ere family \eqref{met:ma-family} and transport
its tail to obtain \eqref{met:endpoint-tail}. Keeping the potentially
concentrating residual measure, we compare on the same sublevel set to
obtain \eqref{met:residual-capacity}. We then take \(k\to\infty\) at
fixed \(t\); the upper Lelong estimate \eqref{met:limit-lelong-upper}
contradicts the lower bound \eqref{met:lelong-lower} as \(t\to0\).

\subsection*{Resolution data and regularized measures}

The log pullback gives
\begin{equation}
 \label{met:adjoint-data}
 L=K_M+\sum_i g_iD_i,
 \qquad g_i<1,
 \qquad g_i<0\Longrightarrow D_i\text{ is }p\text{-exceptional},
\end{equation}
with simple normal crossings support.  Choose a smooth weight \(\ell\)
on \(L\), pulled back from downstairs, and write \(\theta\) for its
curvature.  Fix K\"ahler forms \(\omega\) on \(M\) and \(\omega_Z\) on
\(Z\), and a smooth probability volume \(dV\) on \(M\).  On a singular
space, smooth forms and weights are understood through smooth local
potentials.

For \(0<t\leq1\), put
\begin{equation}
 \label{met:beta-zero}
 \beta_t^0=\theta+t p^*\omega_Z.
\end{equation}
Analytic nefness gives a smooth function \(q_t\leq0\) such that
\(\beta_t^0+\ddc q_t\) is semipositive and positive definite on a dense
open set.  Indeed one can spend only half the available \(t\omega_Z\)
in the nef approximation downstairs.  Consequently
\begin{equation}
 \label{met:exceptional-degree}
 \int_M(\beta_t^0)^n>0,
 \qquad
 \int_{D_i}(\beta_t^0)^{n-1}=0\quad\text{if }g_i<0.
\end{equation}
The second equality follows from the dimension of the image of an
exceptional divisor and the fact that \(\beta_t^0\) is pulled back.

Nefness and weak compactness give a positive current in \([\theta]\).
Let
\begin{equation}
 \label{met:minimal-envelope}
 \Phi=V_\theta
 :=\left(\sup\{v\in\operatorname{PSH}(M,\theta):v\leq0\}\right)^*.
\end{equation}
The weight \(\ell+\Phi\) has minimal singularities.  Suppose, for a
contradiction, that its Lelong number is positive at a point.  In local
coordinates \(x\) centered at that point, choose \(\lambda>0\) such that
\begin{equation}
 \label{met:assumed-pole}
 \Phi(x)\leq\lambda\log|x|^2+O(1).
\end{equation}

In a local frame for \(\OO(D_i)\), write \(s_i\) for its canonical
section and \(d_i=\log|s_i|^2\).  Set \(b=\sum_i g_i d_i\).
The actual identification \eqref{met:adjoint-data} defines a global
volume \(\mu_b\): locally its density is \(e^{\ell-b}\), with the
coordinate volume associated with the canonical frame.  Choose smooth
weights \(d_i^0\) on \(\OO(D_i)\) and set
\begin{equation}
 \label{met:regularized-density}
 d_{i,\varepsilon}
   =\log\bigl(|s_i|^2+\varepsilon e^{d_i^0}\bigr),
 \qquad
 \mu_\varepsilon=e^{\ell-\sum_i g_i d_{i,\varepsilon}},
 \qquad 0<\varepsilon\leq1.
\end{equation}
These local expressions respect the line transitions.  The measures
\(\mu_\varepsilon\) are smooth and positive, and their densities converge
almost everywhere to that of \(\mu_b\).  The simple normal crossings
description and \(g_i<1\) give constants \(p>1\) and \(C\) such that
\begin{equation}
 \label{met:density-moments}
 \left\|\frac{\mu_\varepsilon}{dV}\right\|_{L^p(dV)}\leq C,
 \qquad
 \int_M\left(\frac{\mu_b}{dV}\right)^{-s/(1-s)}dV<\infty
\end{equation}
for some fixed \(0<s<1\).  To obtain the second assertion, take \(s\)
small enough for the finitely many negative coefficients in
\eqref{met:adjoint-data}.  Fix
\begin{equation}
 \label{met:C-zero}
 C_0>\frac{n}{s\lambda}.
\end{equation}

Let \(t\downarrow0\) along a sequence, and for each \(t\) take
\(k=1,2,\ldots\).  We shall choose
\[
 0<\delta=\delta_{t,k}\leq\min(t,1/k),
 \qquad
 \beta=\beta_t^0+\delta\omega,
 \qquad
 V=\int_M\beta^n.
\]
The class of \(\beta\) is K\"ahler, although the chosen smooth
representative \(\beta\) need not be positive.  After choosing
\(\delta\), choose a smooth \(\beta\)-psh function \(\psi=\psi_{t,k}\)
such that
\begin{equation}
 \label{met:psi-approximation}
 \begin{split}
 q_t-k-1&\leq\psi\leq1,\\
 \bigl\|\psi-\max(\Phi,q_t-k)\bigr\|_{L^1(dV)}
    &\leq\min(t,1/k).
 \end{split}
\end{equation}
The maximum is bounded and \(\beta_t^0\)-psh.  Regularization with
arbitrarily small curvature loss gives smooth approximants because
the bounded potential has no positive Lelong numbers
\cite{DemaillyRegularization}.  A smooth convex regularized maximum with
\(q_t-k\), and the upper bound from Hartogs' lemma, give the stated
pointwise bounds.  Finally choose
\(0<\varepsilon=\varepsilon_{t,k}\leq\min(t,1/k)\).
Additional smallness requirements on \(\delta\) and \(\varepsilon\)
will be imposed below, in this order.  None of these data depends on
the parameter \(c\).

For \(0\leq c\leq C_0\), solve
\begin{equation}
 \label{met:ma-family}
 T=\beta+\ddc u>0,
 \qquad
 \rho=\frac{T^n}{V}
   =e^{(1+c)u-c\psi}\mu_\varepsilon
   =e^{u+cH}\mu_\varepsilon,
 \qquad H=u-\psi.
\end{equation}
Here and below \(\rho\) denotes a probability measure.  The
Aubin--Yau theorem applies in the K\"ahler class
\cite{AubinMongeAmpere,YauCalabi}.  More precisely, the positive
coefficient \(1+c\) is the negative-\(\lambda\) case of
\cite[Theorem 7.14]{AubinBook1998}; it gives uniqueness and an
invertible linearized operator.  Thus \(u\) depends smoothly on \(c\).

\subsection*{A capacity estimate independent of the class volume}

We first establish estimates uniform in \(t,k,\delta,\varepsilon,c\)
whenever the displayed conditions hold.  All the representatives
\(\beta\) have a common upper bound by a fixed multiple of \(\omega\).
We use the following standard uniform integrability consequence of
Skoda's theorem and semicontinuity of complex singularity exponents:
for fixed \(B\), there are \(a_B,C_B>0\) such that
\begin{equation}
 \label{met:uniform-skoda}
 \sup_M v=0,\quad \ddc v\geq-B\omega
 \quad\Longrightarrow\quad
 \int_M e^{-a_Bv}dV\leq C_B.
\end{equation}
One obtains uniformity by compactness of normalized quasi-psh
functions, local Skoda integrability at each limit, and
Demailly--Koll\'ar semicontinuity
\cite{SkodaIntegrability,DemaillyKollar}.

Put \(E=V_\beta\), with the same envelope convention as in
\eqref{met:minimal-envelope}.  For each of our K\"ahler classes, \(E\)
is bounded and \(\sup E=0\).  Define the normalized capacity by
\begin{equation}
 \label{met:capacity}
 \operatorname{cap}_\beta(A)
 =\frac1V\sup_{\substack{v\in\operatorname{PSH}(M,\beta)\\E-1\leq v\leq E}}
       \int_A(\beta+\ddc v)^n.
\end{equation}
We suppress the subscript when there is no ambiguity.  In particular,
\(0\leq\operatorname{cap}(A)\leq1\).

There are constants \(a,C>0\), independent of the class volume \(V\),
such that
\begin{equation}
 \label{met:exponential-capacity}
 \mu_\varepsilon(A)
 \leq C\exp\bigl(-a\operatorname{cap}(A)^{-1/n}\bigr)
\end{equation}
for every Borel set \(A\), with the usual value zero on the right when
the capacity is zero.  Here is the normalization argument.  For a
compact nonpluripolar set \(A\), let \(F_A\) be the upper regularization
of the global \(\beta\)-psh extremal with obstacle zero on \(A\), and
put \(m=\sup_M F_A\).  The compact K\"ahler extremal theory gives a
bounded function, at most zero on \(A\) outside a pluripolar set,
maximal outside \(A\), with Monge--Amp\`ere mass \(V\) carried by \(A\)
\cite[Sections 5--7]{GuedjZeriahiCapacity}.

These statements also hold for our possibly nonpositive smooth
representative.  Indeed the defining family has a bounded competitor.
If its suprema were unbounded, sup-normalization, quasi-psh compactness
and a rapidly convergent weighted sum would produce a quasi-psh pole
on \(A\), contradicting nonpluripolarity.  The obstacle constraint
survives upper regularization outside a pluripolar set.  A
quasi-everywhere constraint gives the same envelope: mix a competitor
with arbitrarily small weight on a nonpositive \(\beta\)-psh function
having a pole on the exceptional pluripolar set, whose existence is
the global pluripolarity theorem in a K\"ahler class.  Finally local
balayage off \(A\) proves maximality there.  This explains why no
normalization \(E=0\) is being assumed.

By its defining property and sup-normalization,
\[
 E\leq F_A\leq E+m.
\]
If \(b_A=\max(m,1)\), then
\[
 E+\frac{F_A-E}{b_A}-1
\]
is admissible in \eqref{met:capacity}.  Positivity of mixed products
and the mass support of \(F_A\) give
\begin{equation}
 \label{met:volume-cancellation}
 \operatorname{cap}(A)
 \geq\frac{1}{b_A^nV}\int_A(\beta+\ddc F_A)^n
 =b_A^{-n}.
\end{equation}
In particular \(b_A\geq\operatorname{cap}(A)^{-1/n}\).  The volume
has canceled exactly.  The function \(F_A-m\) satisfies
\eqref{met:uniform-skoda}, and is at most \(-m\) on \(A\) almost
everywhere.  Hence \(dV(A)\leq C e^{-a m}\).  If the capacity is less
than one, \eqref{met:volume-cancellation} implies
\(m\geq\operatorname{cap}(A)^{-1/n}\); capacity one is absorbed in the
constant.  H\"older and \eqref{met:density-moments} prove
\eqref{met:exponential-capacity}.  Pluripolar compact sets have zero
volume, and inner approximation proves the Borel-set assertion.

\subsection*{Initial comparison and a uniformly small shell}

First,
\begin{equation}
 \label{met:upper-u}
 \sup_M u\leq C.
\end{equation}
Otherwise normalize \(u\) by its supremum and take an almost-everywhere
convergent subsequence using quasi-psh compactness.  The limits are
finite almost everywhere.  The measures \(\mu_\varepsilon\) also
subconverge almost everywhere to a density positive almost everywhere,
including when \(\varepsilon\to0\).  Since \(\psi\leq1\), Fatou's
lemma applied to \eqref{met:ma-family} would make its total mass tend
to infinity.  This contradicts \(\int_M\rho=1\).

Fix
\[
 \frac{C_0}{1+C_0}<r<1,
 \qquad v_*=r\psi+(1-r)E.
\]
For a capacity test \(v\) and \(0<\tau\leq1-r\), put
\[
 v_\tau=r\psi+(1-r-\tau)E+\tau v,
 \qquad A_l=\{u<v_*-l\}.
\]
Since \(v_*-\tau\leq v_\tau\leq v_*\), the set
\(B=\{u<v_\tau-l\}\) satisfies
\(A_{l+\tau}\subset B\subset A_l\).  On \(B\), for \(l\geq0\),
\begin{align*}
 (1+c)u-c\psi
 &\leq\bigl((1+c)r-c\bigr)\psi
       +(1+c)(1-r)E-(1+c)l\\
 &\leq C-l.
\end{align*}
We used \((1+c)r-c>0\), \(\psi\leq1\), and \(E\leq0\).
The Bedford--Taylor comparison principle on this same set \(B\)
and mixed-product positivity therefore give
\begin{equation}
 \label{met:first-comparison}
 \tau^n\operatorname{cap}(A_{l+\tau})
 \leq C e^{-l}\mu_\varepsilon(A_l).
\end{equation}

For completeness, write \(g(l)=\operatorname{cap}(A_l)^{1/n}\).
A fixed \(\tau>0\) in \eqref{met:first-comparison} first makes
\(g(l)\) uniformly small for large \(l\).  Then
\eqref{met:exponential-capacity} gives, after enlarging a uniform
constant,
\[
 \tau g(l+\tau)\leq C_1 g(l)^2.
\]
Choose an initial \(l\) so that \(2C_1g(l)\leq1-r\), and successively
take \(\tau=2C_1g(l)\).  Each step halves \(g\), and the sum of the
increments is bounded by a geometric series.  Monotonicity gives
zero capacity above the limiting level.  The resulting almost
everywhere inequality extends to the quasi-psh representatives.
Thus
\begin{equation}
 \label{met:first-lower-bound}
 u\geq r\psi+(1-r)E-C.
\end{equation}
Since \(E\geq\psi-1\), and \(\psi\) has a uniform lower \(L^1\)
bound by \eqref{met:psi-approximation}, we obtain
\begin{equation}
 \label{met:H-normalization}
 H\geq-C,
 \qquad -C\leq\sup_M u\leq C.
\end{equation}
In particular, \(u\) and \(H\) have two-sided pointwise bounds
depending on \(t,k\), but independent of the extra smallness of
\(\delta,\varepsilon\), and independent of the permissible choice of
\(\psi\).

In the rest of the proof, \(o(1)\) means a quantity tending to zero
as \(t\to0\), uniformly in \(k\) and \(c\), after the choices specified
below.  We can choose one fixed large \(R\) for which
\begin{equation}
 \label{met:small-volume-tail}
 dV\{H>R-1\}=o(1).
\end{equation}
To see the uniformity, consider any sequence with \(t\to0\), allowing
both \(k\) and \(c\) to vary.  A subsequence of \(u\) converges in
\(L^1\) and almost everywhere to a \(\theta\)-psh function bounded
above by \(\Phi+C\).  Moreover,
\[
 \| (\Phi-\psi)_+\|_{L^1(dV)}\leq t.
\]
The asserted tail estimate follows by taking \(R>C+2\).
Absolute continuity for \(\mu_b\), the uniform \(L^p\) bound for
\(\mu_\varepsilon\), and \eqref{met:ma-family} now imply that the
shell
\begin{equation}
 \label{met:shell}
 \mathcal S=\{R\leq H\leq R+1\}
\end{equation}
has \(o(1)\) mass for \(\mu_b,\mu_\varepsilon,\rho\).
At \(c=0\), the whole tail \(\{H>R\}\) has \(\rho\)-mass \(o(1)\).
At a general \(c\), only the bounded-shell assertion has so far
been proved.

\subsection*{Differentiating the Monge--Amp\`ere equation}

Set
\[
 h=-\partial_cu,
 \qquad \Delta_T=\tr_T\ddc.
\]
Differentiating both sides of \eqref{met:ma-family} gives
\begin{equation}
 \label{met:differentiated-ma}
 \partial_c\rho
   =\bigl(H-(1+c)h\bigr)\rho
   =-\Delta_T h\,\rho,
 \qquad
 \Delta_T h=(1+c)h-H.
\end{equation}
Also
\[
 \Delta_T H
   =n-\tr_T(\beta+\ddc\psi)\leq n.
\]
At a minimum of \(h\), the last equation in
\eqref{met:differentiated-ma} and \(H\geq-C\) give \(h\geq-C\).
If \(F=(1+c)h-H\), then
\(\Delta_TF\geq(1+c)F-n\); the maximum principle gives
\begin{equation}
 \label{met:h-bounds}
 h\geq-C,
 \qquad
 (1+c)h\leq H+\frac{n}{1+c}.
\end{equation}
Consequently
\begin{equation}
 \label{met:parameter-variation}
 |\partial_c\rho|\leq C(1+H_+)\rho.
\end{equation}
For every fixed \(a>0\), integration by parts on the compact manifold
gives
\begin{equation}
 \label{met:h-energy}
 a\int_M e^{-ah}|\partial h|_T^2\rho
 =\int_M e^{-ah}\Delta_T h\,\rho
 \leq n\int_Me^{-ah}\rho
 \leq C_a.
\end{equation}
Choose a fixed \(C_*\) so large that \(h+C_*>0\), and put
\begin{equation}
 \label{met:transport-form}
 G=(h+C_*)\frac{nT^{n-1}}V.
 \qquad
 \ddc G=-\partial_c\rho.
\end{equation}
The last identity is spatial: at a fixed value of \(c\), \(T\) is
closed, so no parameter derivative of \(T^{n-1}\) enters it.

\subsection*{The order of the smallness choices}

Set
\[
 M_i=\max(0,\lceil-g_i\rceil),
 \quad N=\sum_i M_iD_i,
 \quad d_{N,\varepsilon}=\sum_i M_i d_{i,\varepsilon},
\]
and let \(s_N\) be the canonical section of \(N\).  We impose
\begin{equation}
 \label{met:exceptional-errors}
 \begin{split}
 \sup_M(h+C_*)\int_M
       |\tr_T(\ddc d_{N,\varepsilon})|\,\rho&=o(1),\\
 \int_M\bigl(1-|s_N|_{d_{N,\varepsilon}}^2\bigr)
       |\partial_c\rho|&=o(1).
 \end{split}
\end{equation}
We verify that these choices respect the order already prescribed.

For fixed \(t,k\), Equations \eqref{met:psi-approximation},
\eqref{met:upper-u}, and \eqref{met:h-bounds} give a bound
\(B_{t,k}\) for \(\sup(h+C_*)\) before choosing \(\delta\),
\(\psi\), or \(\varepsilon\).  The signed trace integral is
\begin{equation}
 \label{met:signed-trace}
 \int_M\tr_T(\ddc d_{N,\varepsilon})\rho
   =\frac{n(2\pi)c_1(N)[\beta]^{n-1}}V.
\end{equation}
At \(\delta=0\), its numerator is zero by
\eqref{met:exceptional-degree}, and its denominator is positive for
this fixed \(t\).  Choose \(\delta\), still at most \(\min(t,1/k)\),
so that \(B_{t,k}\) times the modulus of
\eqref{met:signed-trace} is at most \(t\).  Now choose \(\psi\) as in
\eqref{met:psi-approximation}.

It remains to control an absolute trace, not just its signed integral.
The regularization satisfies
\begin{equation}
 \label{met:negative-curvature-tube}
 \ddc d_{i,\varepsilon}
 \geq-C\frac{\varepsilon}
 {|s_i|^2e^{-d_i^0}+\varepsilon}\,\omega.
\end{equation}
For fixed \(t,k,\delta\), the measures
\(T^{n-1}\wedge\omega\) are controlled uniformly in
\(\varepsilon,c\) by a multiple of ordinary \(\omega\)-capacity.
Here is a direct way to make the dependence explicit.  We have
\(\beta\leq A\omega\) and \(|u|\leq B\), for constants fixed at this
stage.  Take \(L_0\geq\max(2A,2B,1)\), enlarging it if necessary,
and set \(v=(u-B)/L_0\).  Then \(-1\leq v\leq0\),
\(\omega+\ddc v\geq\omega/2\), and
\(T\leq L_0(\omega+\ddc v)\).  Thus
\[
 T^{n-1}\wedge\omega
 \leq2L_0^{n-1}(\omega+\ddc v)^n.
\]
The capacity of tubes shrinking to a divisor tends to zero.  The
factor on the right of \eqref{met:negative-curvature-tube} is uniformly
bounded and tends uniformly to zero outside any fixed such tube.
Consequently the negative part of the trace integral tends to zero
uniformly in \(c\) as \(\varepsilon\to0\), at this fixed
\(t,k,\delta\).  The elementary inequality
\[
 \int|f|\,\rho
 \leq\left|\int f\,\rho\right|+2\int f_-\,\rho
\]
then proves the first assertion of \eqref{met:exceptional-errors},
with an error \(O(t)\), after decreasing \(\varepsilon\).

For the second assertion, observe that
\(0\leq|s_N|_{d_{N,\varepsilon}}^2\leq1\), with convergence to one
almost everywhere.  At this stage
\(|\partial_c\rho|\leq C_{t,k}\mu_\varepsilon\).
Uniform integrability from \eqref{met:density-moments} gives the
assertion, again with an error at most \(t\) after decreasing
\(\varepsilon\).  We have therefore chosen \(\delta\), then
\(\psi\), then \(\varepsilon\), for every \(t,k\); all errors are
\(O(t)\), and all choices are independent of \(c\).

\subsection*{A Bochner estimate with a residual term}

Choose a smooth nondecreasing cutoff \(\chi\) with
\[
 \chi(H)=0\text{ for }H\leq R,
 \qquad \chi(H)=1\text{ for }H\geq R+1.
\]
Choose a smooth convex function \(f\) on \(\R\) and a constant \(d\)
such that
\[
 C_0+2\leq f'\leq d-1,
 \qquad f''>0\text{ on }[-R-1,-R].
\]
Define
\begin{equation}
 \label{met:transport-weight}
 W_0=e^{-f(-H)-cH}.
\end{equation}
It is bounded above on \(\mathcal S\).  Since \(H\geq-C\), the lower
bound for \(f'\) and \eqref{met:h-bounds} give, with fixed positive
constants,
\begin{equation}
 \label{met:weight-domination}
 W_0\geq C^{-1}e^{2H},
 \qquad
 W_0\geq C^{-1}(1+H_+^2),
 \qquad
 W_0\geq C^{-1}(h+C_*),
 \qquad
 W_0\geq C^{-1}e^h.
\end{equation}

On \(M^\circ=M\setminus\bigcup_iD_i\), use the following weight on
\(-K_M+N\):
\begin{equation}
 \label{met:bochner-weight}
 P=-\ell-u+b+\sum_iM_i d_i+f(-H).
\end{equation}
The divisor weights are pluriharmonic on this open set.  Therefore
\begin{align}
 \label{met:bochner-curvature}
 \ddc P
 &=-T+(\beta-\theta)-f'(-H)\ddc H
       +f''(-H)\ii\partial H\wedge\bar\partial H\notag\\
 &\geq-dT+f''(-H)\ii\partial H\wedge\bar\partial H.
\end{align}
We used \(\beta-\theta=t p^*\omega_Z+\delta\omega\geq0\) and
\(\ddc H=T-(\beta+\ddc\psi)\leq T\).

We claim that there is an \(N\)-valued section
\(w=s_N\chi(H)-U\) on \(M^\circ\) such that
\begin{equation}
 \label{met:bochner-result}
 \int_{M^\circ}
 \left(|U|_{d_{N,\varepsilon}}^2
       +|\bar\partial w|_{d_{N,\varepsilon},T}^2\right)
 W_0\,\rho=o(1).
\end{equation}
The residual derivative in this assertion is essential: the curvature
lower bound in \eqref{met:bochner-curvature} is not semipositive.

Identify \(N\)-valued sections with top-degree forms valued in
\(-K_M+N\).  Choose a complete K\"ahler form \(\Omega\) on
\(M^\circ\), using logarithmic cusp terms along the simple normal
crossings divisor, and use the complete metrics
\(T_\kappa=T+\kappa\Omega\), \(\kappa>0\).  Let
\[
 g=\bar\partial(s_N\chi(H))
   =s_N\chi'(H)\bar\partial H.
\]
This is a closed square-integrable \((n,1)\)-form for fixed parameters.
For a \(\bar\partial\)-closed \((n,1)\)-form \(\zeta\) in the domain
of \(\bar\partial^*\), complete-metric Bochner--Kodaira and
\eqref{met:bochner-curvature} yield
\begin{equation}
 \label{met:closed-test-estimate}
 |\langle g,\zeta\rangle|^2
 \leq C\left(\int_{\mathcal S}e^{u-f(-H)}\mu_b\right)
       \bigl(\|\bar\partial^*\zeta\|^2+d\|\zeta\|^2\bigr).
\end{equation}
For clarity, on \((n,1)\)-forms the curvature contribution of
\(-dT\) is bounded below by \(-d\|\zeta\|^2\), because
\(T\leq T_\kappa\).  The positive rank-one curvature term controls
contraction with \(\bar\partial H\).  Cauchy--Schwarz in that
direction gives \eqref{met:closed-test-estimate}, since
\((\chi')^2/f''(-H)\) is bounded on the shell.  The squared top-form
density of \(s_N\) in the weight \(P\) is exactly
\(e^{u-f(-H)}\mu_b\).  This density, and hence the estimate's
constant, is independent of the complete base metric.  Approximation
on a complete K\"ahler manifold justifies the indicated weak tests;
see \cite[Sections 4--5]{DemaillyL2}.

The integral in \eqref{met:closed-test-estimate} is \(o(1)\) by the
shell estimate.  Apply the Riesz representation theorem to the
functional defined on pairs
\[
 \bigl(\bar\partial^*\zeta,\sqrt d\,\zeta\bigr)
 \longmapsto\langle g,\zeta\rangle.
\]
After extending this bounded functional, and projecting its second
component to the closed subspace \(\ker\bar\partial\), we obtain
\begin{equation}
 \label{met:residual-riesz}
 \bar\partial U+\sqrt d\,V_1=g,
 \qquad \bar\partial V_1=0,
 \qquad \|U\|^2+\|V_1\|^2=o(1).
\end{equation}
Although the initial tests were closed, this is an equation against
all appropriate tests.  Indeed, orthogonal projection of a test
\(\zeta\in\operatorname{Dom}\bar\partial^*\) to
\(\ker\bar\partial\) preserves its adjoint domain and adjoint value:
the range of \(\bar\partial\) from \((n,0)\)-forms is contained in
that kernel.  It also preserves its pairings with the closed forms
\(g,V_1\).

Let \(\kappa\downarrow0\).  The norms of top-degree \((n,0)\)-forms
are unchanged by the base metric, while the \((n,1)\)-norms increase
to the norm for \(T\).  For a countable decreasing sequence of
\(\kappa\)'s, weak compactness in each preceding norm and a diagonal
subsequence give limits satisfying \eqref{met:residual-riesz} for
\(T\).  Monotone convergence of the norms preserves its bound.

Finally, the ratio of the \(P\)-density for sections to the density
\(|\,\cdot\,|_{d_{N,\varepsilon}}^2W_0\rho\) is
\begin{equation}
 \label{met:density-ratio}
 \prod_i\left(
   \frac{|s_i|^2+\varepsilon e^{d_i^0}}{|s_i|^2}
          \right)^{g_i+M_i}\geq1.
\end{equation}
The same comparison applies to derivative norms measured with \(T\).
Since \(\bar\partial w=\sqrt d\,V_1\), this proves
\eqref{met:bochner-result}.

For fixed parameters, its smooth positive weights imply ordinary
\(L^2\) bounds on \(w\) and its displayed derivative up to the
divisor.  To extend the derivative distributionally, use simple
normal crossings cutoffs whose derivatives are \(O(\tau^{-1})\) on
tubes of volume \(O(\tau^2)\).  Their derivative \(L^2\) norms are
bounded.  The boundary error is bounded by this fixed bound times
the \(L^2\) norm of \(w\) on the shrinking tubes, and tends to zero
by absolute continuity.  Thus \(\bar\partial w\) extends as its
distributional derivative across the divisor.  Ellipticity on
sections gives
\begin{equation}
 \label{met:sobolev-extension}
 w\in W^{1,2}(M,N).
\end{equation}

\subsection*{The nonholomorphic logarithm and one-sided transport}

Write \(d_N=d_{N,\varepsilon}\) for this subsection and put
\(z=1+|w|_{d_N}^2\).  For a smooth section, set
\[
 A_w=\frac{\langle\bar\partial w,w\rangle_{d_N}}z,
\]
where the Hermitian pairing is linear in the first variable.  In a
holomorphic frame normal for the line metric at the point, direct
differentiation gives
\begin{equation}
 \label{met:log-identity}
 \begin{split}
 \ddc\log z={}&\ii\partial A_w
       -\ii\bar\partial\overline{A_w}
       -\frac{|w|_{d_N}^2}{z}\ddc d_N\\
 &+\frac{\ii}{z^2}
       \bigl(\partial w\wedge\bar\partial\bar w
                 -\partial\bar w\wedge\bar\partial w\bigr).
 \end{split}
\end{equation}
The last line is evaluated in that normal frame, or equivalently
using the Chern derivatives.  Its first summand is nonnegative; its
negative summand is controlled by \(|\bar\partial w|^2\).  Thus no
estimate for the positive \(\partial w\) term is required.

Integrate \eqref{met:log-identity} against the positive form \(G\)
in \eqref{met:transport-form}.  By
\eqref{met:weight-domination} and \eqref{met:bochner-result}, the
negative-gradient contribution is at most
\[
 C\int_M(h+C_*)|\bar\partial w|_{d_N,T}^2\rho=o(1).
\]
The curvature cost is bounded by the first error in
\eqref{met:exceptional-errors}.  For the two exact terms, integration
by parts differentiates only \(h\), since \(T\) is closed.  Using
\(|w|/(1+|w|^2)\leq1/2\), their absolute value is at most
\begin{equation}
 \label{met:log-exact-errors}
 C\left(
   \int_M|\bar\partial w|_{d_N,T}^2W_0\rho
   \int_M|\partial h|_T^2W_0^{-1}\rho
       \right)^{1/2}=o(1).
\end{equation}
The second integral is bounded by \eqref{met:h-energy} and
\(W_0^{-1}\leq C e^{-h}\).  Approximation by smooth sections in
\(W^{1,2}\) justifies these calculations for
\eqref{met:sobolev-extension}: the logarithm has bounded first and
second derivatives as a function of the section, and the derivative
products converge in \(L^1\) for fixed smooth data.  We conclude
\begin{equation}
 \label{met:log-integral}
 \int_M\ddc\log z\wedge G\geq-o(1),
 \qquad
 \int_M\log z\,\partial_c\rho\leq o(1).
\end{equation}
The second assertion follows from \eqref{met:transport-form} by
integration by parts.

We can replace \(\log z\) in the last integral by
\(\log(1+\chi(H)^2)\).  First the globally Lipschitz radial function
\(\xi\mapsto\log(1+|\xi|^2)\), followed by Cauchy--Schwarz, gives
\begin{align*}
 &\int_M\left|
 \log(1+|w|_{d_N}^2)
 -\log(1+|s_N\chi(H)|_{d_N}^2)\right|\,|\partial_c\rho|\\
 &\hspace{2em}\leq
 C\int_M|U|_{d_N}(1+H_+)\rho=o(1).
\end{align*}
Here \(W_0\geq C^{-1}(1+H_+^2)\) pays for the second
Cauchy--Schwarz factor.  Next,
\[
 \left|\log(1+|s_N|_{d_N}^2\chi^2)-\log(1+\chi^2)\right|
 \leq1-|s_N|_{d_N}^2,
\]
so the second error in \eqref{met:exceptional-errors} pays for this
replacement.

Set
\[
 \eta(H)=\frac{\log(1+\chi(H)^2)}{\log2}.
\]
The preceding estimate, \(\partial_cH=-h\leq C\), and the shell
support of the bounded nonnegative function \(\eta'\) give
\begin{equation}
 \label{met:one-sided-transport}
 \partial_c\int_M\eta(H)\rho
 =\int_M\eta(H)\partial_c\rho
       +\int_M\eta'(H)(-h)\rho
 \leq o(1).
\end{equation}
The section \(w\) need not be chosen smoothly in \(c\): its estimate
was used separately at each \(c\), whereas the left side here
depends only on the smooth family \eqref{met:ma-family}.
At \(c=0\) the integral is \(o(1)\) by the initial tail estimate.
Integrating \eqref{met:one-sided-transport} up to \(C_0\) proves,
uniformly in \(k\),
\begin{equation}
 \label{met:endpoint-tail}
 \int_{\{H>R+1\}}\rho=o(1)
 \qquad(c=C_0).
\end{equation}

\subsection*{Canceling the residual measure on the comparison set}

Work now at \(c=C_0\).  Take a smooth cutoff \(\xi(H)\) which is zero
for \(H\leq R+1\), one for \(H\geq R+2\), and lies between zero
and one.  The smooth positive measure
\(\xi(H)\rho+t\,dV\) has mass \(a_{t,k}>0\), with
\(a_{t,k}\leq m_t\to0\) uniformly in \(k\), by
\eqref{met:endpoint-tail}.  Solve the prescribed-volume equation
\cite{YauCalabi}
\[
 \frac{(\beta+\ddc v)^n}{V}
     =\frac{\xi(H)\rho+t\,dV}{a_{t,k}},
 \qquad \sup_Mv=0.
\]
The complementary part of \(\rho\) has \(H\leq R+2\).  Therefore
\begin{equation}
 \label{met:residual-domination}
 \rho\leq C_R\mu_\varepsilon
       +m_t\frac{(\beta+\ddc v)^n}{V}.
\end{equation}
Choose \(r_t\to0\) with \(r_t^n\geq m_t\).  We may assume
\(r_t\leq1/2\) for all the remaining \(t\)'s.

We prove a uniform estimate
\begin{equation}
 \label{met:residual-lower-bound}
 u\geq r_tv+(1-r_t)E-C.
\end{equation}
Put \(v_*=r_tv+(1-r_t)E\), and for an arbitrary capacity test
\(\varphi\) and \(0<\tau\leq1-r_t\), put
\[
 v_\tau=r_tv+(1-r_t-\tau)E+\tau\varphi,
 \qquad A_l=\{u<v_*-l\},
 \qquad B=\{u<v_\tau-l\}.
\]
Again \(A_{l+\tau}\subset B\subset A_l\).  On precisely this
same set \(B\), comparison and positivity give
\begin{align*}
 &r_t^n\int_B\frac{(\beta+\ddc v)^n}{V}
 +\tau^n\int_B\frac{(\beta+\ddc\varphi)^n}{V}\\
 &\qquad\leq\int_B\frac{(\beta+\ddc v_\tau)^n}{V}
 \leq\int_B\rho\\
 &\qquad\leq C_R\mu_\varepsilon(B)
       +m_t\int_B\frac{(\beta+\ddc v)^n}{V}.
\end{align*}
Since \(r_t^n\geq m_t\), the residual measure cancels on \(B\).
Taking the supremum over \(\varphi\) yields
\begin{equation}
 \label{met:residual-capacity}
 \tau^n\operatorname{cap}(A_{l+\tau})
 \leq C_R\mu_\varepsilon(A_l).
\end{equation}
There is no comparison of residual masses on different sets.

To start the iteration uniformly, note that \(v_*\leq0\), so
\(A_l\subset\{u<-l\}\).  Equations
\eqref{met:uniform-skoda} and \eqref{met:H-normalization}, followed
by H\"older with \eqref{met:density-moments}, make
\(\mu_\varepsilon(A_l)\) uniformly small as \(l\to\infty\).
A fixed \(\tau=1/4\) in \eqref{met:residual-capacity} therefore
makes its capacity uniformly small.  Combining with
\eqref{met:exponential-capacity}, the same quadratic iteration used
after \eqref{met:first-comparison} applies; the allowed upper bound
\(1-r_t\) is at least \(1/2\).  It proves
\eqref{met:residual-lower-bound} with a constant independent of
\(t,k\).

\subsection*{The fixed-\texorpdfstring{$t$}{t} limit and Lelong contradiction}

Fix one sufficiently small \(t>0\), and let \(k\to\infty\).
Quasi-psh compactness, \eqref{met:H-normalization}, and \(\sup v=0\)
give a subsequence converging in \(L^1\) and almost everywhere to
\(\beta_t^0\)-psh functions \(u_t,v_t\).  Since \(E\geq q_t\),
Equation \eqref{met:residual-lower-bound} gives
\begin{equation}
 \label{met:limit-lelong-upper}
 u_t\geq r_tv_t-C_t,
 \qquad \nu(u_t,x)\leq r_t\nu(v_t,x)\leq C r_t.
\end{equation}
The last constant is independent of \(t\): the normalized potentials
\(v_t\) have a common lower curvature bound, so their Lelong numbers
are uniformly bounded, for example by \eqref{met:uniform-skoda}.
The additive constant \(C_t\) is allowed to depend on \(t\) and
does not affect Lelong numbers.

At fixed \(t\), Equation \eqref{met:psi-approximation} gives
\(\psi_{t,k}\to\Phi\) in \(L^1\).  Choose the subsequence also to
converge almost everywhere.  Since \(\varepsilon_{t,k}\to0\),
Fatou's lemma in \eqref{met:ma-family} at \(c=C_0\) gives
\begin{equation}
 \label{met:fatou-endpoint}
 \int_M e^{(1+C_0)u_t-C_0\Phi}\mu_b\leq1.
\end{equation}
Writing \(f_b=\mu_b/dV\), H\"older and the second moment in
\eqref{met:density-moments} yield
\begin{align*}
 \int_M e^{s(1+C_0)u_t-sC_0\Phi}dV
 &\leq
 \left(\int_M e^{(1+C_0)u_t-C_0\Phi}f_b\,dV\right)^s
 \left(\int_M f_b^{-s/(1-s)}dV\right)^{1-s}\!<\infty.
\end{align*}
By \eqref{met:assumed-pole}, this implies near the chosen point
\begin{equation}
 \label{met:weighted-integrability}
 \int e^{a u_t(x)}|x|^{-2b_0}dV(x)<\infty,
 \qquad a=s(1+C_0),\quad b_0=sC_0\lambda>n.
\end{equation}
Add a local smooth potential for \(\beta_t^0\) to make \(u_t\)
plurisubharmonic.  Its exponential is subharmonic, and the smooth
addition changes the following estimates only by bounded factors.
For \(x\neq0\) sufficiently close to zero, the submean inequality
on \(B(x,|x|/2)\) gives
\begin{align*}
 e^{a u_t(x)}
 &\leq C_t|x|^{-2n}
       \int_{B(x,|x|/2)} e^{a u_t(y)}dV(y)\\
 &\leq C'_t|x|^{2b_0-2n},
\end{align*}
where the second inequality uses
\eqref{met:weighted-integrability} and \(|y|\asymp|x|\) on that
ball.  Hence
\begin{equation}
 \label{met:lelong-lower}
 u_t(x)\leq
 \frac{sC_0\lambda-n}{s(1+C_0)}\log|x|^2+O_t(1).
\end{equation}
The coefficient is positive and independent of \(t\), contradicting
\eqref{met:limit-lelong-upper} as \(t\to0\).  Thus \(\Phi\) has zero
Lelong numbers everywhere.  Any two metrics with minimal
singularities differ by bounded weights, so the conclusion holds
for every such metric on \(L\).
\end{proof}
